\scanpage{152}
\mspage{97}
$p:Y'\longrightarrow Y$ est $\in W$, et de $\eqsim$ pour $q:X'\longrightarrow X$, et $f\in W$ implique \unsure{encore|aussi|donc} $f'\in W$.
\figure{\begin{tikzpicture}[x=1cm,y=1cm]\draw[line width=.7pt] (0,.18)--(.15,0);\end{tikzpicture}}

Revenons \unsure{à|:} l'hyp. p. 95. On aura donc un foncteur
\[
(\underline{\mathrm{Fib}}\,S,W_S=W_S^{\mathcal L})\longrightarrow(\deleted{\unsure{C}}\mathcal M_{/S},W_S^{\mathcal L})
\]
d'où pour les localisés un foncteur horizontal
\diagram{\begin{tikzcd}[ampersand replacement=\&]
(\underline{\mathrm{Fib}}\,S)W_S^{-1}\arrow[r,"i_0"]\arrow[dd,"\approx"] \& (\mathcal M_{/S})(W_S^{\mathcal L})^{-1}\arrow[d,equals,"\mathrm{def}"]\arrow[ddl]\\
{} \& \mathrm{Ho}^{\mathcal L}_W(S)\\
\mathrm{Ho}^{\mathrm{lc}}_W(S)\overset{\mathrm{def}}{=}(\mathcal M_{/S})W_S^{-1} \& {}
\end{tikzcd}}
qui factorise le foncteur vertical induit par l'inclusion $(\underline{\mathrm{Fib}}\,\deleted{\unsure{S}}S,W_S)\longrightarrow(\deleted{\mathrm{Cat}}\mathcal M_{/S},W_S)$ --- foncteur dont on a vu (dans les deux cas envisagés) que c'est une équivalence.
En inversant ce dernier, \inserted{et le composant avec $i_0$,} on trouve donc un foncteur
\margin{\unsure{*}}
\diagram{\begin{tikzcd}[ampersand replacement=\&]
\mathrm{Ho}^{\mathrm{lc}}_W(S)\arrow[r,"i"]\arrow[d,equals,"\mathrm{def}"'] \& \mathrm{Ho}^{\mathcal L}_W(S)\arrow[d,equals,"\mathrm{def}"]\\
(\mathcal M_{/S})W_S^{-1} \& (\mathcal M_{/S})\deleted{\unsure{(}}(W_S^{\mathcal L})^{-1}\text{,}
\end{tikzcd}}
tel que $\rho i\simeq\mathrm{id}$, donc une \uline{rétraction} de $\mathrm{Ho}^{\mathcal L}_W(S)$ sur $\mathrm{Ho}^{\mathrm{lc}}_W(S)$. On s'attend que $i$ soit pl fidèle, et plus précisément encore,
% archive sheet number: 151
